\section{The game and formal model}
\label{sec:game}

\subsection{Informal rules}

Hexo is a two-player placement game on an unbounded hexagonal grid.  The
opening player places one stone at the origin.  Thereafter the players
alternate turns of two placements.  Stones are permanent.  A placement is
legal on any empty cell within hex distance 8 of a stone of either player.
The first player to occupy six consecutive cells on one of the three grid
axes wins.  A win is checked after each placement, so a win on the first
placement ends the game before the second placement of that turn.

We use attacker and defender as fixed roles, denoted by $\A$ and $\D$.  A
ply is one placement, while a normal turn contains two plies.  Everything
below concerns positions after the Opening placement (the origin is occupied;
budgets follow the turn structure of \cref{def:D4}).

\subsection{Formal model}

The first definition fixes the infinite board and the metric used by both
the line geometry and the legality rule.

\begin{definition}[Board, cells, and distance]
\label{def:D1}
The cells are the integer axial lattice $\Cells\ni x=(q,r)$.  For
$x=(q,r)$ and $y=(q',r')$, set
\[
  \dist(x,y)=\max\bigl\{|q-q'|,\ |r-r'|,
                  |(q-q')+(r-r')|\bigr\}.
\]
\end{definition}

The winning patterns are finite segments along the three line directions of
the axial lattice.

\begin{definition}[Axes and windows]
\label{def:D2}
The axis vectors are
\[
  \mathfrak a\in\{(1,0),(0,1),(1,-1)\}.
\]
A \emph{window} is a set
\[
  W=\{s+i\mathfrak a:0\le i\le5\}
\]
for a start cell $s$ and an axis vector $\mathfrak a$.  Every window has
exactly six cells, and every cell lies in exactly 18 windows, from three axes
and six possible offsets on each axis.
\end{definition}

The next fact records the metric consequence of belonging to one window.

\begin{fact}[Window collinearity]
\label{fact:F1}
Two distinct cells of a window differ by a nonzero multiple of the window's
axis vector.  Thus any two cells contained in a common window are collinear
along that axis and have distance at most $5$.
\end{fact}

Figure~\ref{fig:coordinates} shows the coordinate convention and one distance
calculation.

\begin{figure}[tbp]
  \centering
  \begin{tikzpicture}[scale=.88]
    \HexBoardRange{0}{0}{2}
    \foreach \q/\r/\shade in {0/0/22,1/0/10,2/-1/22}{
      \path[fill=black!\shade,draw=none] \HexAt{\q}{\r}
        ++(30:.94) \foreach \angle in {150,210,270,330,30}
        {-- ++(\angle:.94)} -- cycle;
    }
    \node[font=\scriptsize] at \HexAt{-1}{0} {$(-1,0)$};
    \node[font=\scriptsize] at \HexAt{0}{0} {$(0,0)$};
    \node[font=\scriptsize] at \HexAt{1}{0} {$(1,0)$};
    \node[font=\scriptsize] at \HexAt{2}{0} {$(2,0)$};
    \node[font=\scriptsize] at \HexAt{0}{1} {$(0,1)$};
    \node[font=\scriptsize] at \HexAt{2}{-1} {$(2,-1)$};
  \end{tikzpicture}
  \caption{Pointy-top axial coordinates.  The shaded cells show a shortest
  path from $(0,0)$ to $(2,-1)$, and
  $d=\max(|2|,|-1|,|1|)=2$.}
  \label{fig:coordinates}
\end{figure}

A position records the finite board state and the exact placement remaining
in the current turn.

\begin{definition}[Positions]
\label{def:D3}
A position is $P=(\sigma,m,b)$, where $\sigma$ is a finite partial map
$\Cells\rightharpoonup\{\A,\D\}$, the mover is
$m\in\{\A,\D\}$, and $b\in\{1,2\}$ is the mover's number of remaining
placements in the current turn.  Stones are permanent, and
$\Stones(P)=\operatorname{dom}\sigma$.  For $X\in\{\A,\D\}$, let $\bar X$
denote the other player.  For such an $X$ and a window
$W$, define
\[
  \cnt_X(W,P)=|\{c\in W:\sigma(c)=X\}|.
\]
When the position is clear from context, we suppress it and write
$\cnt_X(W)$.
\end{definition}

Legality is local to the existing stones, while a transition must also
account for immediate termination and the two-placement turn structure.

\begin{definition}[Legality and transitions]
\label{def:D4}
A cell $c$ is \emph{legal} in $P$ if and only if
\[
  c\notin\operatorname{dom}\sigma
  \quad\text{and}\quad
  \exists s\in\operatorname{dom}\sigma:\dist(c,s)\le8.
\]
Write $\Legal(P)$ for the set of cells legal in $P$.
The mover places on a legal cell $c$, producing
$\sigma'=\sigma\cup\{c\mapsto m\}$.  If some window $W$ then satisfies
$\cnt_m(W)=6$, the game terminates immediately and $m$ wins.  Wins are
therefore checked after each placement and may occur on the first placement
of a two-placement turn.  If there is no win, the next position is
$ (\sigma',m,1)$ when $b=2$, and $(\sigma',\bar m,2)$ when $b=1$.
\end{definition}

Window states distinguish usable winning lines, blocked lines, and lines that
can be completed within one normal turn.

\begin{definition}[Window states]
\label{def:D5}
A window $W$ is \emph{alive for $X$} if
$\cnt_{\bar X}(W)=0$.  An all-empty window is alive for both players.  It is
\emph{dead} if $\cnt_\A(W)\ge1$ and $\cnt_\D(W)\ge1$.  It is an
\emph{$X$-threat} if it is alive for $X$ and $\cnt_X(W)\ge4$.
\end{definition}

The three panels in \cref{fig:window-states} give the basic mask patterns used
by this terminology.

\begin{figure}[tbp]
  \centering
  \begin{subfigure}[t]{.30\textwidth}
    \centering
    \begin{tikzpicture}[scale=.38]
      \HexPatch{0/0,1/0,2/0,3/0,4/0,5/0}
      \HexWindow[window one]{0}{0}{1}{0}
      \HexAttacker{1}{0}\HexAttacker{4}{0}
    \end{tikzpicture}
    \caption{Alive for $\A$.}
  \end{subfigure}\hfill
  \begin{subfigure}[t]{.30\textwidth}
    \centering
    \begin{tikzpicture}[scale=.38]
      \HexPatch{0/0,1/0,2/0,3/0,4/0,5/0}
      \HexAttacker{1}{0}\HexDefender{4}{0}
    \end{tikzpicture}
    \caption{Dead.}
  \end{subfigure}\hfill
  \begin{subfigure}[t]{.30\textwidth}
    \centering
    \begin{tikzpicture}[scale=.38]
      \HexPatch{0/0,1/0,2/0,3/0,4/0,5/0}
      \HexWindow[window two]{0}{0}{1}{0}
      \HexAttacker{0}{0}\HexAttacker{1}{0}
      \HexAttacker{3}{0}\HexAttacker{4}{0}
    \end{tikzpicture}
    \caption{An $\A$-threat with count $4$.}
  \end{subfigure}
  \caption{Alive, dead, and threat window states.  Filled disks are attacker
  stones and open disks are defender stones.}
  \label{fig:window-states}
\end{figure}

The implemented one-ply analysis tests whether the mover can win in the
remaining placements or whether the opponent's threats exceed the defensive
budget.

\begin{definition}[$\lambda^1$ predicates]
\label{def:D6}
For mover $m$ with budget $b$, $\ownwinnow(P)$ holds if and only if there is
an $m$-alive window $W$ such that
\[
  \cnt_m(W)=5,
  \qquad\text{or}\qquad
  \cnt_m(W)=4\ \text{ and }\ b=2.
\]
The value $\hit(P)$ is the multiset of empty-cell sets of the
$\bar m$-threat windows.  The value $\mhs(P)$ is the minimum number of cells
meeting every set in $\hit(P)$, computed exhaustively for values at most $2$.
The $\lambda^1$ verdict is \textsc{win} when $\ownwinnow(P)$ holds,
\textsc{loss} when it does not hold and $\mhs(P)>b$, and undecided otherwise.
The post-opening soundness of these implemented verdicts is assumed here.
\end{definition}

Absolute placement indices allow later certificates to state an exact win
horizon without treating a two-placement turn as simultaneous.

\begin{definition}[Plies and horizon]
\label{def:D7}
A \emph{ply} is one placement, and ply indices are absolute along a play.  For
a position $P$ with mover $m$, the \emph{defender-placement count before ply
$T$}, written $\DefCount(P,T)$, is the number of plies strictly before $T$ at
which $\D$ places, as determined by $(m,b)$ and alternation.
\end{definition}

The final convention fixes the direction needed when defender replies are
removed from an attacker-win certificate.

\begin{definition}[Dominance direction]
\label{def:D8}
Outcomes are evaluated for the attacker $\A$, with
$\mathrm{WIN}\succ\nonWIN$.  A \emph{sound dismissal} of a defender reply
$c$ at a node $N$ means that, if the certificate proves that $\A$ wins against
all searched replies, then $\A$ also wins against $c$.
\end{definition}

\begin{figure}[tbp]
  \centering
  \begin{tikzpicture}[scale=.46]
    \HexPatch{0/-4,1/-4,-1/-3,0/-3,1/-3,-1/-2,0/-2,1/-2,2/-2,
      -1/-1,0/-1,1/-1,2/-1,3/-1,4/-1,5/-1,
      -2/0,-1/0,0/0,1/0,2/0,3/0,4/0,5/0,
      -2/1,-1/1,0/1,1/1,2/1,3/1,4/1,
      -3/2,-2/2,-1/2,0/2,1/2,
      -4/3,-3/3,-2/3,-1/3,0/3,
      -5/4,-4/4,-3/4,-5/5,-4/5}
    \HexWindow[fill=black!8,draw=none]{-1}{0}{1}{0}
    \HexWindow[fill=black!18,draw=none]{0}{-3}{0}{1}
    \HexWindow[fill=black!30,draw=none]{-4}{4}{1}{-1}
    \HexMark{0}{0}{dot}{}
    \draw[-{Stealth[length=4pt]},line width=.55pt]
      \HexAt{5.75}{0} -- \HexAt{6.30}{0};
    \node[font=\scriptsize,anchor=west] at \HexAt{6.45}{0} {$(1,0)$ axis};
    \draw[-{Stealth[length=4pt]},line width=.55pt]
      \HexAt{0}{3.70} -- \HexAt{0}{4.25};
    \node[font=\scriptsize,anchor=south west] at \HexAt{0}{4.40}
      {$(0,1)$ axis};
    \draw[-{Stealth[length=4pt]},line width=.55pt]
      \HexAt{2.70}{-2.70} -- \HexAt{3.25}{-3.25};
    \node[font=\scriptsize,anchor=north west] at \HexAt{3.40}{-3.40}
      {$(1,-1)$ axis};
  \end{tikzpicture}
  \caption{One six-cell window on each axis through the central cell, shown
  with three fill shades.  Sliding a window so that the cell occupies offsets
  $0$ through $5$ gives $3$ axes $\times$ $6$ offsets $=18$ windows through
  the marked cell.}
  \label{fig:axes-windows}
\end{figure}
